% \textbf{First paragraph:}
% \begin{itemize}
%     \item Context: Diffusion MRI provides quantitative microstructural biomarkers
%     \item Increasing use of ML for prediction (age, diagnosis, cognition, etc.)
%     \item Lack of standardized benchmark
% \end{itemize}
Diffusion magnetic resonance imaging (dMRI) enables the in vivo quantification of tissue 
microstructure and provides biologically meaningful biomarkers sensitive to axonal organization, 
myelination, and tissue complexity \cite{novikov2019quantifying}. Leveraging these microstructural 
descriptors for predictive modeling has become increasingly common, with applications spanning 
demographic estimation \cite{chen2024deep,he2022model}, clinical diagnosis 
\cite{egle2022prediction,wen2021reproducible}, and cognitive characterization 
\cite{kotikalapudi2025replicability}. The combination of diffusion-derived features and machine 
learning is therefore often viewed as a powerful framework for data-driven neurobiological inference.

% \noindent \textbf{Second paragraph: literature review:}
% \begin{itemize}
%     \item  Literature review of techniques: Variability of the proposed approach (ML-based, deep-learning based etc)
%     \item Problems: High variability in reported performance across studies + Limited reproducibility 
% \end{itemize}

Yet despite rapid methodological development and steadily increasing reported accuracies, a fundamental 
question remains unresolved: how robust are these gains? Reported performance varies widely across 
studies, tasks, and datasets, and improvements are frequently attributed to architectural innovation or 
feature engineering. However, direct comparison between studies is rarely possible, and it remains 
unclear whether observed gains reflect genuine microstructural signal, dataset-specific properties, or 
sensitivity to methodological choices. In the absence of standardized evaluation, diffusion MRI 
prediction risks becoming a comparison of configurations rather than a cumulative scientific endeavor.

% \noindent \textbf{Third paragraph: reasons for these problems}
% \begin{itemize}
%     \item Reason 1: Heterogeneous preprocessing pipelines
%     \item Reason 2: Inconsistent feature representations (voxel-wise vs ROI-based), not analysis cross microstructure features
%     \item Reason 3: Evaluation protocols not directly comparable + results on different datasets
%     \item Reason 4: Missing code
%     \item Our solution: "To address these challenges, we propose: 1) an open-source and standardized benchmark; 2) of various machine learning techniques in dMRI; 3) on several tasks (classification, regression) 4) encompassing N datasets 
% \end{itemize}

Several recurrent sources of inconsistency contribute to this fragmentation. Preprocessing pipelines 
and quality-control procedures differ substantially across cohorts. Microstructural information is 
represented through heterogeneous strategies (e.g., voxel-wise maps, ROI summaries, or learned 
embeddings), yet evaluations rarely disentangle representation effects from modeling effects. 
Evaluation protocols vary in split definitions, cross-validation schemes, and hyperparameter selection, 
often conflating methodological flexibility with performance gains. Finally, incomplete release of code 
and configuration details further limits reproducibility. Together, these factors make it difficult to 
determine which reported improvements are stable, generalizable, and biologically meaningful.


% \noindent \textbf{Fourth paragraph: your contributions}
% \begin{itemize}
%     \item Standardized preprocessing pipeline for diffusion MRI
%     \item Implemented various representation of the microstructural input data
%     \item Unified feature extraction strategy
%     \item Controlled evaluation protocol (train/test splits)
%     \item Comparison of multiple feature representations
%     \item Integration of both classical and DL predictors
%     \item Public release of code and benchmark configuration
%     \item End-to-end analysis
% \end{itemize}
% In this work, we introduce an open-source, end-to-end benchmark for predictive modeling with 
% diffusion MRI that targets reproducibility and fair comparison across studies. Our contributions 
% are: (i) a standardized preprocessing and data preparation pipeline; (ii) a unified feature-extraction 
% framework supporting multiple microstructural representations for gray and white matter; (iii) a 
% controlled evaluation protocol with consistent, leakage-free train/validation/test splits and 
% cross-validation; (iv) a systematic comparison of feature representations and predictors, spanning 
% classical machine learning baselines and deep-learning models; and (v) the public release of code, 
% configurations, and scripts to facilitate transparent reuse and extension.

To address these limitations, we introduce a standardized, open-source diffusion MRI benchmark that 
enables controlled, leakage-free comparison across datasets, tissue compartments, microstructural 
representations, and model families. Rather than focusing solely on peak accuracy, the benchmark 
explicitly quantifies performance variability across folds and configurations, isolates 
representation and tissue-specific effects, and evaluates the incremental value of diffusion-derived 
features beyond non-diffusion baselines. By unifying preprocessing, feature extraction, model 
selection, and evaluation under a single experimental framework, this work aims to shift diffusion MRI 
prediction from isolated performance claims toward variance-aware, reproducible benchmarking.

Our contributions are fourfold: (i) a standardized preprocessing and data preparation pipeline; (ii) a 
unified feature-extraction framework supporting multiple microstructural representations for gray and 
white matter; (iii) a consistent evaluation protocol with leakage-free cross-validation and 
configuration-level statistical analysis; and (iv) the public release of code, splits, and experimental 
configurations to enable transparent reuse and extension. Collectively, this benchmark provides the 
foundation for disentangling genuine microstructural signal from configuration-driven variability and 
for establishing comparable evaluation standards in diffusion MRI prediction.
